\documentclass[11pt,a4paper]{article}
\usepackage[utf8]{inputenc}
\usepackage[T1]{fontenc}
\usepackage[french]{babel}
\usepackage{amsmath,amssymb}
\usepackage{booktabs}
\usepackage{geometry}
\geometry{margin=2.5cm}

\title{Implementation et Evaluation de Methodes ICA Classiques et Stochastiques}
\author{Abdoulaye Diallo}
\date{21 mars 2026}

\begin{document}
\maketitle
\tableofcontents
\newpage

\section{Introduction}

\subsection{Contexte}
L'Analyse en Composantes Independantes (ICA) vise a separer des sources latentes statistiquement independantes a partir d'observations melangees lineairement. Les methodes classiques comme Infomax \cite{bell1995information} et FastICA \cite{hyvarinen1999fast} sont des references solides, mais les variantes stochastiques sont importantes pour des jeux de donnees volumineux et pour le traitement en flux.

Dans ce projet on cherche a comparer des methodes ICA en termes de qualite de separation et de comportement numerique. La metrique centrale retenue est l'indice d'Amari \cite{amari1996new}, qui quantifie l'ecart entre la matrice de performance estimee et une structure permutation-dilatation ideale.
On propose egalement une contribution creative en explorant l'impact de la gaussianite des sources sur les performances ICA, en utilisant une gaussienne generalisee avec un parametre de kurtosis variable.
\subsection{Structure du projet}
Le projet se construit de la maniere suivante pour atteindre:
\begin{itemize}
\item implementation de plusieurs methodes ICA, incluant des variantes stochastiques,
\item validation de l'implementation via des tests synthetiques,
\item comparaison systematique des methodes sur des jeux de donnees synthetiques avec evaluation via l'indice d'Amari,
\item Et enfin etude de l'impact de la gaussianite sur les performances ICA.
\end{itemize}

\section{Fondements Theoriques}
On commence par les bases de l'ICA, en particulier le modele lineaire, l'hypothese de non-gaussianite et le critere Infomax.
\subsection{Le modele ICA}
Le modele lineaire est:
\begin{equation}
x = A s,
\end{equation}
ou $s$ est un vecteur de sources independantes et $A$ la matrice de melange. On cherche une matrice de demelange $V$ telle que $y = Vx$ recupere les sources a permutation et echelle pres.

\subsection{Hypothese de non-gaussianite}
L'identifiabilite de l'ICA repose sur la non-gaussianite des sources (au plus une source gaussienne). Cette hypothese motive l'usage de contrastes non lineaires (par exemple logcosh/tanh) dans les algorithmes pratiques.

\subsection{Critere Infomax}
Pour un echantillon $x$ et $y=Vx$, la log-vraisemblance echantillon est ecrite:
\begin{equation}
\ell(V;x)=\log|\det(V)|+\sum_{j=1}^{D}\log p_j(y_j).
\end{equation}
Avec la fonction score $g(u)=\frac{d}{du}\log p(u)$, le gradient echantillon est:
\begin{equation}
\nabla_V\ell(V;x)=V^{-T}+g(y)x^T.
\end{equation}
Apres whitening, on impose en pratique $VV^T\approx I$, ce qui conduit a l'approximation
$V^{-T}\approx V$ et a une mise a jour de type
$V \leftarrow V + \eta\left(V+\mathbb{E}[g(Vx)x^T]\right)$.


\section{Methodes}

\subsection{Algorithmes implementes}

\subsubsection{Infomax avec gradient batch}
Le critere optimise est la log-vraisemblance Infomax. Dans le cadre du projet
(donnees centrees puis blanchies), la mise a jour est appliquee en ascension de gradient
avec mini-batchs et re-orthogonalisation numerique de $V$ a chaque iteration.

\subsubsection{FastICA}
Pour FastICA, l'implementation \texttt{sklearn.decomposition.FastICA} est utilisee (mode parallel, non linearite logcosh) avec calcul de l'indice d'Amari.
\subsubsection{Choix de la fonction score}
Pour rester coherent avec l'enonce, on utilise la definition score stricte
$g(u)=\frac{d}{du}\log p(u)$. Avec une prior logcosh/super-gaussienne,
on retient $g(u)=-\tanh(u)$.

\subsubsection{Hyperparametres}
Les parametres utilises pour FastICA sont:
\begin{itemize}
\item algorithm = parallel,
\item fun = logcosh,
\item max\_iter = 1000,
\item tol = 1e-6,
\item random\_state = 42.
\end{itemize}



\subsubsection{SGD-ICA}
SGD-ICA implemente une approximation mini-batch du gradient Infomax:
\begin{equation}
V \leftarrow V + \eta\left(V+\frac{1}{m}\sum_{i=1}^{m} g(Vx_i)x_i^T\right).
\end{equation}
La contrainte d'orthogonalite apres whitening est imposee numeriquement
par re-orthogonalisation (QR).

\subsubsection{Adam-ICA}
Adam-ICA utilise exactement le meme gradient que SGD-ICA, mais remplace le pas
constant par les moments adaptes d'Adam \cite{kingma2014adam}. Cela permet de
tester si une dynamique adaptative ameliore la stabilite ou la vitesse de convergence.

\section{Experiences}

\subsection{Protocole experimental}
Le protocole experimental repose sur des tests synthetiques de separation:
\begin{itemize}
\item generation de sources non gaussiennes independantes,
\item melange lineaire par une matrice $A$ inversible,
\item pretraitement (centrage + whitening),
\item estimation de $\hat V$,
\item evaluation via $C=\hat V A$ et indice d'Amari.
\end{itemize}
Deux campagnes principales ont ete menees:
\begin{itemize}
\item scalabilite en dimension $D$ (avec $N$ fixe),
\item scalabilite en taille d'echantillon $N$ (avec $D$ fixe).
\end{itemize}

\subsection{Datasets}

\subsubsection{Donnees synthetiques}
Les experiences utilisent des melanges synthetiques, avec repetition sur plusieurs seeds.
Le cas de base est $D=3$, $N=5000$. Pour les etudes de scalabilite, $D$ ou $N$ est ensuite
balaye sur une grille de valeurs.



\subsection{Metriques d'evaluation}
L'evaluation se fait principalement via l'indice d'Amari, qui mesure la proximite de la matrice de performance $C=\hat V A$ a une matrice de permutation-dilatation ideale. Un indice proche de 0 indique une separation quasi-parfaite, tandis que des valeurs plus elevees indiquent des erreurs de separation.
\subsubsection{Indice d'Amari}
L'indice d'Amari est calcule a partir de $C=\hat V A$:
\begin{equation}
r(C)=\frac{1}{2D(D-1)}\left[\sum_i\left(\frac{\sum_j|c_{ij}|}{\max_k|c_{ik}|}-1\right)+\sum_j\left(\frac{\sum_i|c_{ij}|}{\max_k|c_{kj}|}-1\right)\right].
\end{equation}
Une valeur proche de 0 indique une bonne separation.

\subsection{Resultats}

\subsubsection{Comparaison des algorithmes}
Cas de base ($D=3$, $N=5000$):

\begin{table}[h]
\centering
\begin{tabular}{lccc}
\toprule
Algorithme & Amari Index & Temps (s) & Convergence \\
\midrule
FastICA (sklearn) & 0.0147 & 0.01 & Oui \\
SGD-ICA (Infomax) & 0.1324 & 0.92 & Oui \\
Adam-ICA (Infomax) & 0.2553 & 0.92 & Oui \\
\bottomrule
\end{tabular}
\caption{Comparaison sur le cas de base (seed fixe)}
\end{table}

Scalabilite en dimension ($N=5000$, moyennes sur 3 seeds):
\begin{table}[h]
\centering
\begin{tabular}{lcc}
\toprule
Algorithme & Amari moyen (D=3\,$\to$\,40) & Temps moyen (D=3\,$\to$\,40) \\
\midrule
FastICA & 0.0124\,$\to$\,0.0129 & 0.014s\,$\to$\,0.584s \\
SGD-ICA & 0.1516\,$\to$\,0.3044 & 0.521s\,$\to$\,14.315s \\
Adam-ICA & 0.3073\,$\to$\,0.3116 & 0.488s\,$\to$\,12.151s \\
\bottomrule
\end{tabular}
\caption{Tendance de la scalabilite en dimension}
\end{table}

Scalabilite en nombre d'echantillons ($D=20$, budget stochastique en epochs fixes):
\begin{table}[h]
\centering
\begin{tabular}{lcc}
\toprule
Algorithme & Amari moyen (N=1000\,$\to$\,10000) & Temps moyen (N=1000\,$\to$\,10000) \\
\midrule
FastICA & 0.0289\,$\to$\,0.0089 & 0.067s\,$\to$\,0.470s \\
SGD-ICA & 0.3451\,$\to$\,0.3206 & 0.085s\,$\to$\,9.542s \\
Adam-ICA & 0.3474\,$\to$\,0.3294 & 0.067s\,$\to$\,8.646s \\
\bottomrule
\end{tabular}
\caption{Tendance de la scalabilite en taille d'echantillon}
\end{table}



\section{Rôle de la gaussianite dans l'ICA}

L'apport original retenu est l'etude de l'impact de la gaussianite sur l'ICA via une experience faisant varier le parametre de kurtosis d'une gaussienne generalisee.

\subsection{Methode proposee}


\subsection{Resultats}



\section{Discussion}

\subsection{Analyse des resultats}
Les experiences montrent un comportement robuste de FastICA en qualite de separation
(indice d'Amari tres faible et stable), aussi bien quand $D$ augmente que quand $N$ augmente.
Avec les reglages actuels, SGD-ICA et Adam-ICA convergent, mais restent sur un plateau de
performance sensiblement moins bon.

\subsection{Limites identifiees}
Les resultats stochastiques restent sensibles a la convention de score, aux hyperparametres,
et a la frequence de re-orthogonalisation. Le budget de calcul (iterations, epochs) influence
fortement la comparaison de temps entre methodes batch et stochastiques.


\subsection{Comparaison entre methodes}
Dans l'etat actuel du projet, FastICA domine en precision et en temps. Cela ne contredit pas
l'interet theorique des methodes stochastiques, mais indique qu'un tuning plus fin est
necessaire pour observer un avantage pratique en grande echelle.


\subsection{Insights sur la contribution}
La partie experimentation a permis de clarifier un point important: la formulation mathematique
du gradient (notamment le choix de la score function) doit etre strictement coherente avec
la log-vraisemblance de reference pour obtenir un comportement numerique interpretable.



\section{Conclusion}

\subsection{Resume des apports}
Le projet a produit:
\begin{itemize}
\item une implementation fonctionnelle de FastICA, SGD-ICA et Adam-ICA autour du critere Infomax,
\item une validation sur donnees synthetiques via l'indice d'Amari,
\item deux etudes de scalabilite (en $D$ et en $N$) avec mesures de qualite et de cout.
\end{itemize}

\subsection{Perspectives}
Les prochaines etapes naturelles sont: (i) comparaison a budget temps strictement egal,
(ii) tuning systematique des hyperparametres stochastiques, (iii) reduction du cout de
re-orthogonalisation, et (iv) integration concise des resultats de la partie 2 sur la gaussianite.



\bibliographystyle{plain}
\bibliography{references}

\end{document}
